\section{Cardinals and the Axiom of Choice}\label{Ch1:Sec:Cardinals}

In this section, we discuss how big the levels of the $V$-hierarchy are. The answer is elementary at the bottom of the hierarchy, but by the time we reach $V_{\omega + 1}$ it already depends on the Axiom of Choice. We end by defining two hierarchies of cardinals, only one of which needs choice.

\subsection{The Size of $V_{\omega}$}

We begin with the finite levels.

\begin{boxproposition}
    For all $n < \omega$, the set $V_{n}$ is finite.
\end{boxproposition}

Indeed, their sizes are $0$, $2^0$, $2^{2^{0}}$, and so on.

So $V_{\omega}$ is the first infinite level. It turns out to be countable, and there are two ways of seeing it.

\begin{boxproposition}
    $V_{\omega}$ is countable.
\end{boxproposition}
\begin{proof}[Proof using AC]
    Write
    \begin{align*}
        V_{\omega} = \bigcup_{n < \omega} V_n
    \end{align*}
    AC tells us that a countable union of countable sets is countable.
\end{proof}
\begin{proof}[Proof that avoids AC]
    By recursion on $n < \omega$, we can define $R_n \ssq V_n \times V_n$ so that
    \begin{enumerate}[label = (\arabic*)]
        \item $\parenth{V_n, R_n}$ is a well-ordering
        \item $R_{n+1}$ end-extends $R_n$
    \end{enumerate}
    Then, we'll be able to put $R_{\omega} = \bigcup_{n < \omega} R_n$ and show that $R_{\omega}$ is a well-ordering of type $\omega$.

    Since $\parenth{V_n, R_n}$ is a well-ordering of a finite set, we may enumerate $V_n$ as
    \begin{align*}
        a_0 \, R_n \, a_1 \, R_n \, \cdots \, R_n \, a_{k-1}
    \end{align*}
    where $k = \abs{V_n}$. Every member of $V_{n+1} = \Powset\of{V_n}$ is a subset of $V_n$, so we may write any $x, y \in V_{n+1}$ as
    \begin{align*}
        x = \setst{a_i}{i \in I}
        \qquad \text{and} \qquad
        y = \setst{a_i}{i \in J}
    \end{align*}
    for unique $I, J \ssq k$. Now, put $x \, R_{n+1} \, y$ iff
    \begin{enumerate}[label = (\roman*)]
        \item $x, y \in V_n$ and $x \, R_n \, y$.
        \item $x \in V_n$ and $y \in V_{n + 1} \setminus V_n$.
        \item $x, y \in V_{n+1} \setminus V_n$, and $i \in J$, where $i$ is the largest natural number lying in exactly one of $I$ and $J$.
    \end{enumerate}
    Case (iii) makes sense because $x \neq y$ forces $I \neq J$, so some natural number lies in exactly one of them and there are only finitely many candidates. The relation it defines is the ordering of the subsets of $V_n$ by binary numerals, reading $x$ as the string whose $i$th digit records whether $i \in I$, so it is a linear ordering.

    So $R_{n+1}$ is $R_n$ followed by a linear ordering of the finitely many sets in $V_{n+1} \setminus V_n$, which makes it a well-ordering and gives us (1). It agrees with $R_n$ on $V_n$ and places all of $V_n$ before all of $V_{n+1} \setminus V_n$, so $V_n$ is an initial segment of $\parenth{V_{n+1}, R_{n+1}}$, which gives us (2).

    The point of the construction is that $R_{n+1}$ is defined outright from $R_n$, with no choices made along the way, and the recursion starts from $R_0 = \emptyset$ on $V_0 = \emptyset$. The $R_n$ increase, so $R_{\omega}$ linearly orders $V_{\omega}$. Chaining (2) makes $V_n$ an initial segment of $\parenth{V_m, R_m}$ for every $m \geq n$, so each $x \in V_{\omega}$ has all of its $R_{\omega}$-predecessors inside the finite $V_n$ it came from, and every proper initial segment of $R_{\omega}$ is therefore finite. Since $\omega \ssq V_{\omega}$, the ordering is infinite, so sending each $x$ to its number of predecessors is an isomorphism onto $\omega$.
\end{proof}

\subsection{Cardinality Without Choice}

The reason the second proof is worth having is that without AC, countable unions are badly behaved.

\begin{remark}
    If ZF is consistent, then so is ZF $+$ ``there's a countable union of pairs that has no cardinality''.
\end{remark}

The same problem arises one level further up.

\begin{remark}
    Without AC, we don't know whether the set $V_{\omega + 1} = \Powset\of{V_{\omega}}$ even has a \textit{cardinality}. With AC,
    \begin{align*}
        \abs{V_{\omega + 1}} = \abs{\Powset\of{V_{\omega}}} = 2^{\aleph_0}
    \end{align*}
    where $\aleph_0 = \omega = \abs{V_{\omega}}$.
\end{remark}

\subsection{Cardinal Successors}

It turns out that the successor of a cardinal can still be defined without choice, provided we go about it the right way.

Given AC, we know that $2^{\kappa} > \kappa$ and we can define $\kappa^{+}$ to be the least cardinal strictly greater than $\kappa$. Without AC, we still know, given a cardinal $\kappa$, that
\begin{align*}
% [CORRECTED] was: $\setst{R \ssq \kappa}{\parenth{\kappa, R} \text{ is a well-ordering}}$ --- $R$ has to be a relation on $\kappa$ for $\parenth{\kappa, R}$ to be a structure, and no ordinal is an ordered pair, so for every $\kappa \geq 2$ the displayed class was empty
    \setst{R \ssq \kappa \times \kappa}{\parenth{\kappa, R} \text{ is a well-ordering}}
\end{align*}
is a set. We can still have Mostowski Collapse for well-orderings. We can then conclude that
\begin{align*}
    \setst{\text{types}\of{\kappa, R}}{\parenth{\kappa, R} \text{ is a well-ordering}}
\end{align*}
is a set of ordinals. Then, we can set $\kappa^{+}$ to simply be the supremum of this set.

\subsection{The Alephs and the Beths}

We can now define the two standard hierarchies of cardinals, of which only the second needs the Axiom of Choice.

Without AC, define
\begin{align*}
    \aleph_0 &:= \omega \\
    \aleph_{\alpha + 1} &:= \parenth{\aleph_{\alpha}}^{+} \\
    \aleph_{\beta} &:= \sup_{\alpha < \beta}\of{\aleph_{\alpha}} \text{ if } \beta \text{ is a limit ordinal} 
\end{align*}
With AC, define
\begin{align*}
    \beth_0 &:= \omega \\
% [CORRECTED] was: $\beth_{\alpha + 1} := 2^{\beth_0}$ --- as written the successor clause never consults $\beth_{\alpha}$, so every $\beth_{\alpha}$ with $\alpha \geq 1$ collapses to $2^{\aleph_0}$
    \beth_{\alpha + 1} &:= 2^{\beth_{\alpha}} \\
    \beth_{\beta} &:= \sup_{\alpha < \beta}\of{\beth_{\alpha}} \text{ if } \beta \text{ is a limit ordinal}
\end{align*}

% The lecture of 31 Aug ends mid-sentence here. The raw notes read ``It's clear'' and stop;
% most likely the start of a comparison of the two hierarchies, but there is more than one
% thing it could have been, so it is marked as a gap rather than guessed at. Raised with the author.
It's clear \sorry 
